\documentclass[10pt]{article}
\usepackage[paperwidth=170mm,paperheight=225mm,left=17mm,right=17mm,top=14mm,bottom=14mm]{geometry}
\usepackage{amsmath}
\usepackage[mathrm=sym]{unicode-math}
\setmainfont{NewCM10-Book.otf}
\setmathfont{STIXTwoMath-Regular.otf}
\pagestyle{empty}
\setlength{\parindent}{0pt}
\setlength{\parskip}{4.5pt}
\begin{document}
\fontsize{10.6}{14.0}\selectfont
{\fontsize{29}{32}\selectfont Downward Glance}\par
{\fontsize{15}{18}\selectfont Proof of the integer value}\par
\vspace{3pt}
The opening floor contains one integral from $-3$ to $3$ in the real variable $l$.
The large parenthesis ends just before $\mathrm{d}l$. Each continuation begins
with $+$. Reading the printed integrand by ordinary precedence gives 2,507
products in 216 lines.

\smallskip{\fontsize{14}{17}\selectfont Six intervals suffice}\par
Cut the integral at $-2,-1,0,1,2$. On every resulting open unit interval
$\lfloor l\rfloor$ and $\lceil l\rceil$ are constants, and $|l|$ is affine.
The other moduli have affine arguments of constant sign there.
Their zeroes and discrete jumps occur only at interval endpoints.
Endpoint values exist and have no effect on these integrals.

Factorials, binomials and finite sums or products are evaluated directly.
A repeated decimal block of $m$ digits with integer value $b$ contributes
$b/(10^m-1)$. Each printed radical has a nonnegative rational square as its
radicand. A two-by-two determinant contributes $ad-bc$ in row order.
An inner moment with upper limit $B(l)\geq0$ contributes
\[
   \int_0^{B(l)} t^p\,\mathrm{d}t=\frac{B(l)^{p+1}}{p+1},\qquad \int_0^{|l|}t^4\,\mathrm{d}t=\frac{|l|^5}{5}.
\]
The remaining inner integrals have constant bounds and polynomial integrands.
All divisors are nonzero. Consequently each product on each interval is a
polynomial over the rationals. Its degree is at most 13.

\smallskip{\fontsize{14}{17}\selectfont Exact computation}\par
The score evaluator expands these polynomials and integrates each monomial.
A second evaluator parses the complete printed TeX. For degree bound $d$, it
samples the polynomial at the $d+1$ interior points
$j+k/(d+2)$, $1\leq k\leq d+1$, on $(j,j+1)$.
Integrating the Lagrange basis gives rational weights exact through degree $d$.
Both calculations agree separately on every line. For the unrounded integral
$I$ and the integer $N$, the resulting equality is
\[
 I=N+\frac{17602437396215302323978349}{25410946577841365997456000},\qquad 0<I-N<1 .
\]
\smallskip{\fontsize{14}{17}\selectfont The printed floor}\par
\[
   \lfloor I\rfloor=N=-373\,571\,877\,419\,178\,597\,233\,231\,125\,653 .
\]
\end{document}
